% appendix_final.tex — Supplementary material

\section{Experimental Details}
\label{app:experimental_details}

\subsection{Hardware and Software}
All experiments are conducted on NVIDIA A100-80GB GPUs.
We use PyTorch 2.10 with CUDA 12.8, HuggingFace Transformers (v4.44+),
and bfloat16 precision.
Greedy decoding ($\tau{=}0$) is used throughout; no sampling or
nucleus filtering is applied.

\subsection{Budget Control}
Token budgets are controlled via the \texttt{max\_new\_tokens} parameter
in HuggingFace \texttt{model.generate()}, which caps the \emph{total}
number of newly generated tokens.
For thinking mode (\texttt{enable\_thinking=True}), this budget is shared
between the reasoning trace (\texttt{<think>...</think>}) and the final
answer.
For non-thinking mode (\texttt{enable\_thinking=False}), the entire budget
is available for the answer.
This ensures a fair comparison: both modes operate under the same total
output-token budget.

\subsection{Answer Extraction}
We use a multi-level extraction pipeline:
\begin{enumerate}[nosep,leftmargin=*]
    \item Search for \texttt{\textbackslash boxed\{...\}} patterns
    \item Search for \texttt{\#\#\#\#} markers (GSM8K convention)
    \item Search for ``Final answer:'' patterns
    \item Fall back to the last number in the output
\end{enumerate}
When thinking mode exhausts its budget without producing a final answer,
we optionally apply a \emph{projection pass}: a short continuation
(max 16 tokens) to extract a numerical answer from the truncated output.
This is controlled by the \texttt{--projection\_on\_missing\_final} flag.

\subsection{Dataset Details}
\paragraph{GSM8K.} We use the standard test split ($n{=}1{,}319$ problems)~\citep{cobbe2021gsm8k}.
A 200-sample subset (seed 42) is used for the budget sweep in
Table~\ref{tab:thinking-tax}; all other analyses use the complete set
unless otherwise noted.

\paragraph{MATH-500.} We use a 500-problem subset of the MATH benchmark~\citep{hendrycks2021math},
following the split used by~\citet{lightman2024lets}.
Due to computational constraints, we evaluate 40-sample subsets
across multiple data seeds (42, 404, 505, 606, 707) and report averages.

% ============================================================================
\section{Full Data Tables}
\label{app:full-tables}

\subsection{Qwen3-8B: 200-Sample Subset}

Table~\ref{tab:thinking-tax} in the main text presents the complete
200-sample subset results.
All experiments use seed=42, temperature=0 (greedy decoding).
Below we provide the complete data for all tested budgets:

\begin{table}[h]
\centering
\small
\caption{Complete budget sweep for Qwen3-8B on GSM8K 200-sample subset.}
\label{tab:full-8b-200}
\begin{tabular}{llrrrr}
\toprule
\textbf{Mode} & \textbf{Budget} & \textbf{Accuracy} & \textbf{Avg Tokens} & \textbf{Early Stop} & \textbf{Has Final} \\
\midrule
nothink & 32 & 3.0\% & 32 & 0.0\% & 0.0\% \\
nothink & 64 & 12.0\% & 64 & 2.0\% & 0.0\% \\
nothink & 128 & 54.5\% & 111 & 43.5\% & 0.0\% \\
nothink & 256 & 89.0\% & 140 & 92.0\% & 1.5\% \\
nothink & 512 & 94.0\% & 145 & 99.5\% & 1.5\% \\
\midrule
thinking & 128 & 2.0\% & 128 & 0.0\% & 0.0\% \\
thinking & 256 & 22.0\% & 255 & 2.0\% & 0.5\% \\
thinking & 512 & 66.5\% & 442 & 47.5\% & 5.5\% \\
\bottomrule
\end{tabular}
\end{table}

\subsection{Qwen3-8B: Full GSM8K ($n$=1,319)}

\begin{table}[h]
\centering
\small
\caption{Key configurations for Qwen3-8B on full GSM8K ($n{=}1{,}319$).}
\label{tab:full-8b-1319}
\begin{tabular}{llrrrr}
\toprule
\textbf{Mode} & \textbf{Budget} & \textbf{Accuracy} & \textbf{Avg Tokens} & \textbf{Early Stop} & \textbf{Has Final} \\
\midrule
nothink & 128 & 50.8\% & 113 & --- & --- \\
nothink & 256 & 87.5\% & 146 & 88.8\% & 0.6\% \\
\midrule
thinking & 128 & 3.0\% & 128 & --- & --- \\
thinking & 256 & 18.0\% & 255 & --- & --- \\
thinking & 512 & 65.2\% & 460 & 41.3\% & 6.1\% \\
\bottomrule
\end{tabular}
\end{table}

\subsection{Qwen3.5-27B: Full GSM8K ($n$=1,319)}

\begin{table}[h]
\centering
\small
\caption{Qwen3.5-27B thinking mode results on full GSM8K.
The model achieves dramatically lower accuracy than the 8B model at
all budgets due to longer reasoning chains.}
\label{tab:27b-full}
\begin{tabular}{lrrrr}
\toprule
\textbf{Config} & \textbf{Accuracy} & \textbf{Avg Tokens} & \textbf{Has Final} & \textbf{Projection Rate} \\
\midrule
thinking@128 & 3.6\% & 144 & 0.0\% & 100.0\% \\
thinking@256 & 7.9\% & 272 & 0.0\% & 100.0\% \\
thinking@512 & 18.3\% & 528 & 0.7\% & 99.3\% \\
\bottomrule
\end{tabular}
\end{table}

% ============================================================================
\section{DeepSeek-R1 Cross-Model Validation}
\label{app:deepseek}

We validate our findings on DeepSeek-R1-Distill-Llama-8B using MATH-500
with multiple data seeds to ensure robustness.

\begin{table}[h]
\centering
\small
\caption{DeepSeek-R1-Distill-Llama-8B on MATH-500 subsets (40 samples per seed).
The thinking tax and token waste phenomena persist across model families.}
\label{tab:deepseek-math500}
\begin{tabular}{lrrrr}
\toprule
\textbf{Budget} & \textbf{Accuracy} & \textbf{Avg Tokens} & \textbf{Final Rate} & \textbf{Utilization} \\
\midrule
1024 & 22.5\% & 936 & 27.5\% & 91.4\% \\
2048 & 32.5\% & 1644 & 50.0\% & 80.3\% \\
4096 & 47.5\% & 2599 & 72.5\% & 63.5\% \\
\bottomrule
\end{tabular}
\end{table}

Key findings that generalize from GSM8K:
\begin{itemize}[nosep,leftmargin=*]
    \item \textbf{Token utilization collapses with budget}: from 91.4\%
      at $b{=}1024$ to 63.5\% at $b{=}4096$.
    \item \textbf{Natural stop correlates with accuracy}: samples that
      produce a final answer marker achieve consistently higher accuracy
      than truncated samples.
    \item \textbf{Diminishing returns}: accuracy gains from 2048$\to$4096
      (+15.0pp) require 2$\times$ the tokens, yielding lower marginal
      efficiency.
\end{itemize}

% ============================================================================
\section{TOWN Simulation Details}
\label{app:town-simulation}

Our per-sample TOWN simulation uses the \textbf{full GSM8K test set}
($n{=}1{,}319$) where we have matched nothink@256 and thinking@512
results for each question.

\paragraph{Routing statistics.}
\begin{itemize}[nosep,leftmargin=*]
    \item \textbf{Stage 1 accepted}: 1{,}171/1{,}319 (88.8\%) samples stop
      before the $B_1{=}256$ budget. Among these, accuracy is
      \textbf{94.4\%} (1{,}106/1{,}171 correct).
    \item \textbf{Stage 2 routed}: 148/1{,}319 (11.2\%) samples exhaust the
      Stage~1 budget. Among these, thinking@512 achieves \textbf{62.8\%}
      accuracy (93/148 correct).
    \item \textbf{Combined}: 1{,}199/1{,}319 = \textbf{90.9\%} accuracy at
      \textbf{199} average tokens (including Stage~1 probe cost for
      all samples).
\end{itemize}

\paragraph{Error analysis of routed samples.}
Of the 148 samples routed to Stage~2:
\begin{itemize}[nosep,leftmargin=*]
    \item \textbf{64} were incorrectly answered by nothink but correctly recovered by thinking@512 (genuine wins)
    \item \textbf{29} were correct in both modes (correctly answered despite routing)
    \item \textbf{36} were incorrect in both modes (genuinely hard problems)
    \item \textbf{19} were actually correct in nothink@256 (despite hitting budget)
      but thinking@512 answered incorrectly (routing regret)
\end{itemize}
The 19 routing regret cases represent a small inefficiency (1.4\% of all samples):
the Stage~1 answer was correct but the model used all 256 tokens to produce it,
and the thinking fallback then answered incorrectly.
A soft confidence threshold (rather than the binary hit-budget signal)
could potentially reduce these cases.
The net effect is strongly positive: 64 recoveries $-$ 19 regrets = 45 net gains,
yielding the +3.4\,pp improvement over \texttt{nothink@256}.

% ============================================================================
\section{Thinking Efficiency Frontier}
\label{app:frontier}

We categorize GSM8K problems by the minimum thinking budget required
for correct answers (see Table~\ref{tab:efficiency-frontier} in the
main text).
The 31.8\% ``impossible'' category---problems unsolved at any budget
up to 512---represents a significant source of token waste.

\paragraph{Oracle analysis.}
An oracle that perfectly identifies problem difficulty and assigns
the minimum sufficient budget to solvable questions achieves:
\begin{itemize}[nosep,leftmargin=*]
    \item \textbf{68.2\%} accuracy ($+3.0$pp over fixed-512)
    \item \textbf{401} average tokens ($-21.7$\% vs.\ fixed-512)
\end{itemize}
The accuracy gain comes from avoiding ``overthinking'' cases where
extended reasoning causes the model to revise a correct intermediate
answer.

% ============================================================================
\section{Reproducibility}
\label{app:reproducibility}

\paragraph{Models.}
\begin{itemize}[nosep,leftmargin=*]
    \item Qwen3-8B: \texttt{Qwen/Qwen3-8B} from HuggingFace
    \item Qwen3.5-27B: \texttt{Qwen/Qwen3.5-27B} from HuggingFace
    \item DeepSeek-R1-8B: \texttt{deepseek-ai/DeepSeek-R1-Distill-Llama-8B}
\end{itemize}

\paragraph{Random seeds.}
All experiments use seed 42 for reproducibility.
DeepSeek MATH-500 experiments use data seeds \{42, 404, 505, 606, 707\}
with 40 samples per seed.

\paragraph{Compute.}
Total compute budget: approximately 200 A100-GPU-hours across all
experiments reported in this paper.
